% ------------------------------------------------------------
\section{Introduction and motivation}
% ------------------------------------------------------------

Diffusive search in heterogeneous environments is ubiquitous: molecules searching
for binding sites in a crowded cytoplasm~\cite{Elf2005,Metzler2014}, pedestrians
in populated corridors, bees navigating a congested nest toward a food
source~\cite{Peleg2018, Ocko2014}, and protein first passage in cellular
consortia~\cite{Schuss2015,Benichou2010}.  In all of these systems, local crowder
density modulates the particle's effective mobility; the diffusion coefficient
becomes position-dependent, $D=D(x)$, rather than constant.  Even in the
``slowly varying'' regime where the motion is still normal (mean-square
displacement linear in time), a spatially heterogeneous $D(x)$ can substantially
alter search times~\cite{Redner2001}.

The central question of the present work is:
\begin{quote}
\itshape
Given a partially absorbing target at $x=0$ with discoverability $\kappa$, does a
spatially dependent diffusivity $D(x)$ speed up or slow down first passage
relative to the homogeneous case $D\equiv D_0$?
\end{quote}
We derive the relevant Fokker--Planck equation from a microscopic lattice random
walk, establish and solve the backward equation for the MFPT on the ring with
periodic boundary conditions and a $\delta$-type trap, and carry out a complete
case study for the piecewise-linear (``tent'') profile
\begin{equation}
D(x)\;=\;D_0 \;+\; D_1\!\left(1-\frac{|x|}{\pi}\right),
\qquad x\in[-\pi,\pi]. \label{eq:tent}
\end{equation}
The analytic results are checked numerically by three independent methods:
finite-difference on the backward equation, Crank--Nicolson on the forward FPE,
and Monte Carlo simulation of the It\^o SDE.  This numerical component makes
explicit the connection between the PDE formulation, the stochastic-process
formulation, and the particle-level simulation~\cite{Gardiner2009,Pavliotis2014,Kloeden1992}.

The work develops tools from the course (Fokker--Planck equations, adjoint
formulations, boundary conditions, self-adjointness, conservative
finite-difference discretisation, Crank--Nicolson time stepping) in a concrete
setting where the answer is both analytically tractable and biologically
meaningful.
